\documentclass[11pt,a4paper]{article}

\usepackage[utf8]{inputenc}
\usepackage[T1]{fontenc}
\usepackage[french]{babel}
\usepackage[a4paper,margin=2.2cm]{geometry}
\usepackage{amsmath,amssymb,amsthm,mathtools}
\usepackage{lmodern}
\usepackage{enumitem}
\usepackage{hyperref}
\usepackage{xcolor}
\usepackage{fancyhdr}
\usepackage{tikz}

\hypersetup{
    colorlinks=true,
    linkcolor=black,
    urlcolor=blue
}

\setlength{\parindent}{0pt}
\setlength{\parskip}{0.6em}

\pagestyle{fancy}
\fancyhf{}
\fancyhead[L]{Fiche d’exercices --- Algèbre linéaire Chapitre 3}
\fancyhead[R]{Espaces vectoriels et sous-espaces}
\fancyfoot[C]{\thepage}

\begin{document}

\begin{center}
    {\LARGE \textbf{Fiche d’exercices --- Chapitre 3}}\\[0.4em]
    {\Large \textbf{Espaces vectoriels et sous-espaces}}\\[0.6em]
    Algèbre linéaire
\end{center}

\textbf{AIDE de lecture type d’exo :} APP = application directe, METH = méthode, DEM = démonstration, PREP = réflexion plus poussée. Les étoiles représentent un gradient de difficulté de 1 à 5.

\tableofcontents

\newpage

\section*{Exercices d’application directe}
\addcontentsline{toc}{section}{Exercices d’application directe}

\textbf{Exercice 1 [APP $\star$]}\\
\textbf{Reconnaître des objets vectoriels}\\
Parmi les ensembles suivants, dire lesquels sont naturellement des espaces vectoriels réels :
\begin{enumerate}[label=\alph*)]
    \item $\mathbb{R}^2$,
    \item $\mathbb{R}^3$,
    \item $\mathbb{R}_+$,
    \item l’ensemble des polynômes réels,
    \item l’ensemble des fonctions réelles définies sur $\mathbb{R}$.
\end{enumerate}

\textbf{Exercice 2 [APP $\star$]}\\
\textbf{Vecteur nul et opposé}\\
Dans un espace vectoriel $E$, compléter les égalités suivantes :
\begin{enumerate}[label=\alph*)]
    \item $u+ \dots = u$,
    \item $u + (\dots)=0_E$,
    \item $(-1)u = \dots$,
    \item $0\cdot u = \dots$,
    \item $\lambda 0_E = \dots$.
\end{enumerate}

\textbf{Exercice 3 [APP $\star\star$]}\\
\textbf{Stabilité simple}\\
Dire si les parties suivantes de $\mathbb{R}^2$ sont stables par addition, par multiplication par un réel, par aucune des deux :
\begin{enumerate}[label=\alph*)]
    \item $\{(x,y)\in\mathbb{R}^2\mid y=0\}$,
    \item $\{(x,y)\in\mathbb{R}^2\mid x+y=1\}$,
    \item $\{(x,y)\in\mathbb{R}^2\mid x\ge 0\}$,
    \item $\{(x,y)\in\mathbb{R}^2\mid y=2x\}$,
    \item $\{(x,y)\in\mathbb{R}^2\mid xy=0\}$.
\end{enumerate}

\textbf{Exercice 4 [APP $\star\star$]}\\
\textbf{Sous-espace ou non ?}\\
Déterminer si les ensembles suivants sont des sous-espaces vectoriels de $\mathbb{R}^2$ :
\begin{enumerate}[label=\alph*)]
    \item $F=\{(x,y)\in\mathbb{R}^2\mid y=3x\}$,
    \item $G=\{(x,y)\in\mathbb{R}^2\mid y=3x+1\}$,
    \item $H=\{(x,y)\in\mathbb{R}^2\mid x-y=0\}$,
    \item $K=\{(x,y)\in\mathbb{R}^2\mid x^2+y^2=0\}$.
\end{enumerate}

\textbf{Exercice 5 [APP $\star\star$]}\\
\textbf{Sous-espaces de $\mathbb{R}^3$}\\
Déterminer si les ensembles suivants sont des sous-espaces vectoriels de $\mathbb{R}^3$ :
\begin{enumerate}[label=\alph*)]
    \item $\{(x,y,z)\in\mathbb{R}^3\mid x+y+z=0\}$,
    \item $\{(x,y,z)\in\mathbb{R}^3\mid x-y+2z=1\}$,
    \item $\{(x,y,z)\in\mathbb{R}^3\mid z=0\}$,
    \item $\{(x,y,z)\in\mathbb{R}^3\mid x=y=z\}$.
\end{enumerate}

\textbf{Exercice 6 [APP $\star\star$]}\\
\textbf{Combinaisons linéaires dans $\mathbb{R}^2$}\\
Dire si les vecteurs suivants sont des combinaisons linéaires des vecteurs donnés :
\begin{enumerate}[label=\alph*)]
    \item $(3,5)$ de $(1,0)$ et $(0,1)$ ;
    \item $(4,8)$ de $(1,2)$ ;
    \item $(1,1)$ de $(1,-1)$ et $(1,1)$ ;
    \item $(2,1)$ de $(1,1)$ et $(1,-1)$.
\end{enumerate}

\textbf{Exercice 7 [APP $\star\star$]}\\
\textbf{Sous-espace engendré simple}\\
Déterminer explicitement :
\begin{enumerate}[label=\alph*)]
    \item $\mathrm{Vect}((1,0))$,
    \item $\mathrm{Vect}((0,1))$,
    \item $\mathrm{Vect}((1,2))$,
    \item $\mathrm{Vect}((1,0),(0,1))$.
\end{enumerate}

\textbf{Exercice 8 [APP $\star\star$]}\\
\textbf{Exemples sur les fonctions}\\
Dans l’espace vectoriel des fonctions réelles définies sur $\mathbb{R}$, dire si les ensembles suivants sont des sous-espaces vectoriels :
\begin{enumerate}[label=\alph*)]
    \item l’ensemble des fonctions paires,
    \item l’ensemble des fonctions impaires,
    \item l’ensemble des fonctions positives,
    \item l’ensemble des fonctions telles que $f(0)=0$.
\end{enumerate}

\textbf{Exercice 9 [APP $\star\star$]}\\
\textbf{Matrices}\\
Dans $M_2(\mathbb{R})$, déterminer si les ensembles suivants sont des sous-espaces vectoriels :
\begin{enumerate}[label=\alph*)]
    \item les matrices diagonales,
    \item les matrices de trace nulle,
    \item les matrices de déterminant nul,
    \item les matrices symétriques.
\end{enumerate}

\textbf{Exercice 10 [APP $\star\star$]}\\
\textbf{Intersection et somme, premiers exemples}\\
Dans $\mathbb{R}^3$, on considère
\[
F=\{(x,y,z)\in\mathbb{R}^3\mid z=0\},\qquad
G=\{(x,y,z)\in\mathbb{R}^3\mid y=0\}.
\]
Décrire :
\begin{enumerate}[label=\alph*)]
    \item $F\cap G$,
    \item $F+G$.
\end{enumerate}

\newpage

\section*{Exercices de méthode}
\addcontentsline{toc}{section}{Exercices de méthode}

\textbf{Exercice 11 [METH $\star\star$]}\\
\textbf{Utiliser le bon critère}\\
Montrer que
\[
F=\{(x,y)\in\mathbb{R}^2\mid x-2y=0\}
\]
est un sous-espace vectoriel de $\mathbb{R}^2$ en utilisant le critère de stabilité par combinaison linéaire.

\textbf{Exercice 12 [METH $\star\star$]}\\
\textbf{Montrer qu’un ensemble n’est pas un sous-espace}\\
Montrer que
\[
A=\{(x,y)\in\mathbb{R}^2\mid x+y=1\}
\]
n’est pas un sous-espace vectoriel de $\mathbb{R}^2$ de deux façons différentes.

\textbf{Exercice 13 [METH $\star\star\star$]}\\
\textbf{Décrire un sous-espace engendré}\\
Déterminer explicitement
\[
\mathrm{Vect}\big((1,1),(1,-1)\big)\subset \mathbb{R}^2.
\]
On demandera une rédaction par analyse-synthèse.

\textbf{Exercice 14 [METH $\star\star\star$]}\\
\textbf{Paramétrer un sous-espace}\\
Décrire explicitement l’ensemble
\[
F=\{(x,y,z)\in\mathbb{R}^3\mid x+y-z=0\}.
\]
Puis l’écrire comme sous-espace engendré par deux vecteurs bien choisis.

\textbf{Exercice 15 [METH $\star\star\star$]}\\
\textbf{Fonctions et stabilité}\\
Montrer que l’ensemble
\[
E=\{f:\mathbb{R}\to\mathbb{R}\mid f(1)=0\}
\]
est un sous-espace vectoriel de l’espace des fonctions réelles sur $\mathbb{R}$.

\textbf{Exercice 16 [METH $\star\star\star$]}\\
\textbf{Droite affine ou droite vectorielle ?}\\
Parmi les ensembles suivants, distinguer ceux qui sont des droites vectorielles de ceux qui sont seulement des droites affines :
\begin{enumerate}[label=\alph*)]
    \item $\{(x,y)\in\mathbb{R}^2\mid y=2x\}$,
    \item $\{(x,y)\in\mathbb{R}^2\mid y=2x+3\}$,
    \item $\{(x,y)\in\mathbb{R}^2\mid x=0\}$,
    \item $\{(x,y)\in\mathbb{R}^2\mid x=1\}$.
\end{enumerate}

\textbf{Exercice 17 [METH $\star\star\star$]}\\
\textbf{Sous-espace engendré dans $\mathbb{R}^3$}\\
Déterminer explicitement
\[
\mathrm{Vect}\big((1,0,0),(0,1,0)\big)\subset\mathbb{R}^3.
\]
Même question avec
\[
\mathrm{Vect}\big((1,1,0),(1,-1,0)\big).
\]

\textbf{Exercice 18 [METH $\star\star\star$]}\\
\textbf{Montrer qu’un vecteur appartient à un sous-espace engendré}\\
Dire si le vecteur $(2,3,1)$ appartient à
\[
\mathrm{Vect}\big((1,1,0),(0,1,1)\big).
\]
Même question pour $(1,0,1)$.

\textbf{Exercice 19 [METH $\star\star\star$]}\\
\textbf{Intersection de deux sous-espaces}\\
Dans $\mathbb{R}^3$, on considère
\[
F=\{(x,y,z)\in\mathbb{R}^3\mid x+y+z=0\},
\qquad
G=\{(x,y,z)\in\mathbb{R}^3\mid x-y=0\}.
\]
Déterminer explicitement $F\cap G$.

\textbf{Exercice 20 [METH $\star\star\star$]}\\
\textbf{Somme de sous-espaces}\\
Dans $\mathbb{R}^3$, on considère
\[
F=\mathrm{Vect}((1,0,0)),\qquad G=\mathrm{Vect}((0,1,0),(0,0,1)).
\]
Décrire explicitement $F+G$.

\newpage

\section*{Exercices de démonstration}
\addcontentsline{toc}{section}{Exercices de démonstration}

\textbf{Exercice 21 [DEM $\star\star$]}\\
\textbf{Unicité du vecteur nul}\\
Montrer que dans un espace vectoriel, le vecteur nul est unique.

\textbf{Exercice 22 [DEM $\star\star$]}\\
\textbf{Unicité de l’opposé}\\
Montrer que tout vecteur d’un espace vectoriel admet un unique opposé.

\textbf{Exercice 23 [DEM $\star\star$]}\\
\textbf{Calcul usuel}\\
Montrer que pour tout vecteur $u$ d’un espace vectoriel,
\[
0u=0_E.
\]

\textbf{Exercice 24 [DEM $\star\star$]}\\
\textbf{Autre calcul usuel}\\
Montrer que pour tout vecteur $u$ d’un espace vectoriel,
\[
(-1)u=-u.
\]

\textbf{Exercice 25 [DEM $\star\star\star$]}\\
\textbf{Produit nul vectoriel}\\
Montrer que si $\lambda u=0_E$, alors $\lambda=0$ ou $u=0_E$.

\textbf{Exercice 26 [DEM $\star\star\star$]}\\
\textbf{Intersection de sous-espaces}\\
Montrer que l’intersection de deux sous-espaces vectoriels d’un même espace vectoriel est un sous-espace vectoriel.

\textbf{Exercice 27 [DEM $\star\star\star$]}\\
\textbf{Somme de sous-espaces}\\
Montrer que la somme de deux sous-espaces vectoriels est un sous-espace vectoriel.

\textbf{Exercice 28 [DEM $\star\star\star$]}\\
\textbf{Le sous-espace engendré est un sous-espace}\\
Soient $u_1,\dots,u_p$ des vecteurs d’un espace vectoriel $E$. Montrer que
\[
\mathrm{Vect}(u_1,\dots,u_p)
\]
est un sous-espace vectoriel de $E$.

\textbf{Exercice 29 [DEM $\star\star\star$]}\\
\textbf{Plus petit sous-espace contenant une famille}\\
Soit $F$ un sous-espace vectoriel de $E$ contenant des vecteurs $u_1,\dots,u_p$. Montrer que
\[
\mathrm{Vect}(u_1,\dots,u_p)\subset F.
\]

\textbf{Exercice 30 [DEM $\star\star\star$]}\\
\textbf{Union de sous-espaces : attention}\\
Dans $\mathbb{R}^2$, on pose
\[
F=\{(x,0)\mid x\in\mathbb{R}\},\qquad
G=\{(0,y)\mid y\in\mathbb{R}\}.
\]
Montrer que $F$ et $G$ sont des sous-espaces vectoriels, mais que $F\cup G$ n’en est pas un.

\textbf{Exercice 31 [DEM $\star\star\star\star$]}\\
\textbf{Fonctions paires et impaires}\\
Montrer que l’ensemble des fonctions paires est un sous-espace vectoriel de $\mathcal{F}(\mathbb{R},\mathbb{R})$, puis faire de même avec l’ensemble des fonctions impaires.

\textbf{Exercice 32 [DEM $\star\star\star\star$]}\\
\textbf{Un exemple matriciel}\\
Montrer que l’ensemble des matrices symétriques de $M_n(\mathbb{R})$ est un sous-espace vectoriel de $M_n(\mathbb{R})$.

\newpage

\section*{Exercices de réflexion}
\addcontentsline{toc}{section}{Exercices de réflexion}

\textbf{Exercice 33 [PREP $\star\star\star$]}\\
\textbf{Un ensemble trompeur}\\
Dans $\mathbb{R}^2$, on considère
\[
A=\{(x,y)\in\mathbb{R}^2\mid xy=0\}.
\]
Dire si $A$ est un sous-espace vectoriel. On justifiera soigneusement.

\textbf{Exercice 34 [PREP $\star\star\star$]}\\
\textbf{Description explicite}\\
Décrire explicitement
\[
\mathrm{Vect}\big((1,2,0)\big)\subset\mathbb{R}^3.
\]
Quel objet géométrique obtient-on ?

\textbf{Exercice 35 [PREP $\star\star\star$]}\\
\textbf{Plan engendré}\\
Décrire explicitement
\[
\mathrm{Vect}\big((1,0,1),(0,1,1)\big)\subset\mathbb{R}^3.
\]

\textbf{Exercice 36 [PREP $\star\star\star\star$]}\\
\textbf{Trouver une équation cartésienne}\\
On considère
\[
F=\mathrm{Vect}\big((1,1,0),(1,0,1)\big)\subset\mathbb{R}^3.
\]
Décrire $F$ par une condition linéaire portant sur les coordonnées $(x,y,z)$.

\textbf{Exercice 37 [PREP $\star\star\star\star$]}\\
\textbf{Sous-espace défini par deux conditions}\\
Déterminer explicitement l’ensemble
\[
F=\{(x,y,z)\in\mathbb{R}^3\mid x+y+z=0,\ x-y=0\}.
\]
Le décrire ensuite comme sous-espace engendré par un seul vecteur.

\textbf{Exercice 38 [PREP $\star\star\star\star$]}\\
\textbf{Une somme égale à tout l’espace}\\
Montrer que
\[
\mathrm{Vect}((1,1),(1,-1))=\mathbb{R}^2.
\]

\textbf{Exercice 39 [PREP $\star\star\star\star$]}\\
\textbf{Une somme qui n’est pas directe au sens intuitif}\\
Dans $\mathbb{R}^2$, on considère
\[
F=\mathrm{Vect}((1,0)),\qquad G=\mathrm{Vect}((2,0)).
\]
Déterminer $F$, $G$, $F\cap G$ et $F+G$. Que remarque-t-on ?

\textbf{Exercice 40 [PREP $\star\star\star\star$]}\\
\textbf{Fonctions et combinaisons linéaires}\\
Montrer que toute fonction de la forme
\[
f(x)=ax+b
\]
appartient à
\[
\mathrm{Vect}(1,x)
\]
dans l’espace des fonctions de $\mathbb{R}$ dans $\mathbb{R}$.\\
Réciproquement, décrire tous les éléments de $\mathrm{Vect}(1,x)$.

\textbf{Exercice 41 [PREP $\star\star\star\star\star$]}\\
\textbf{Un sous-espace ou une réunion de droites ?}\\
Dans $\mathbb{R}^2$, on considère
\[
A=\{(x,y)\in\mathbb{R}^2\mid y=x \text{ ou } y=-x\}.
\]
Déterminer si $A$ est un sous-espace vectoriel.

\textbf{Exercice 42 [PREP $\star\star\star\star\star$]}\\
\textbf{Sous-espace de polynômes}\\
Dans $\mathbb{R}[X]$, on considère l’ensemble
\[
F=\{P\in\mathbb{R}[X]\mid P(1)=0\}.
\]
Montrer que $F$ est un sous-espace vectoriel.

\textbf{Exercice 43 [PREP $\star\star\star\star\star$]}\\
\textbf{Polynômes pairs}\\
Dans $\mathbb{R}[X]$, on considère
\[
E=\{P\in\mathbb{R}[X]\mid P(-X)=P(X)\}.
\]
Montrer que $E$ est un sous-espace vectoriel de $\mathbb{R}[X]$.

\textbf{Exercice 44 [PREP $\star\star\star\star\star$]}\\
\textbf{Analyse-synthèse un peu plus subtile}\\
Déterminer explicitement
\[
\mathrm{Vect}\big((1,1,0),(1,0,1),(0,1,1)\big)\subset\mathbb{R}^3.
\]

\textbf{Exercice 45 [PREP $\star\star\star\star\star$]}\\
\textbf{Problème de synthèse final}\\
Pour chacune des affirmations suivantes :
\begin{enumerate}[label=\alph*)]
    \item dire si elle est vraie ou fausse ;
    \item si elle est vraie, la démontrer ;
    \item si elle est fausse, donner un contre-exemple.
\end{enumerate}

\begin{enumerate}[label=\arabic*)]
    \item L’union de deux sous-espaces vectoriels d’un même espace vectoriel est toujours un sous-espace vectoriel.
    \item L’intersection de deux sous-espaces vectoriels d’un même espace vectoriel est toujours un sous-espace vectoriel.
    \item Toute partie de $\mathbb{R}^2$ contenant $(0,0)$ est un sous-espace vectoriel de $\mathbb{R}^2$.
    \item Si $u$ et $v$ appartiennent à un sous-espace vectoriel $F$, alors tout vecteur de la forme $\lambda u+\mu v$ appartient à $F$.
\end{enumerate}

\newpage

\section*{Pour aller plus loin}
\addcontentsline{toc}{section}{Pour aller plus loin}

\textbf{Exercice 46 [PREP $\star\star\star\star\star$]}\\
\textbf{Sous-espace défini par une relation de somme}\\
Dans l’espace vectoriel des suites réelles, on considère l’ensemble
\[
F=\{(u_n)_{n\ge 0}\mid \forall n\ge 0,\ u_{n+1}=u_n\}.
\]
Montrer que $F$ est un sous-espace vectoriel.

\textbf{Exercice 47 [PREP $\star\star\star\star\star$]}\\
\textbf{Équation fonctionnelle linéaire simple}\\
Dans l’espace vectoriel des fonctions de $\mathbb{R}$ dans $\mathbb{R}$, montrer que l’ensemble
\[
E=\{f\mid \forall x\in\mathbb{R},\ f(x+1)=f(x)\}
\]
est un sous-espace vectoriel.

\textbf{Exercice 48 [PREP $\star\star\star\star\star$]}\\
\textbf{Vecteurs colinéaires cachés}\\
Déterminer explicitement
\[
\mathrm{Vect}\big((1,2),(2,4),(3,6)\big)\subset\mathbb{R}^2.
\]

\textbf{Exercice 49 [DEM $\star\star\star\star$]}\\
\textbf{Cas d’égalité de deux sous-espaces engendrés}\\
Montrer que
\[
\mathrm{Vect}(u)=\mathrm{Vect}(\lambda u)
\]
pour tout vecteur $u$ et tout scalaire non nul $\lambda$.

\textbf{Exercice 50 [PREP $\star\star\star\star\star$]}\\
\textbf{Synthèse générale}\\
Dans $\mathbb{R}^3$, on considère
\[
F=\{(x,y,z)\in\mathbb{R}^3\mid x+y+z=0\},
\qquad
G=\mathrm{Vect}((1,-1,0)).
\]
\begin{enumerate}[label=\alph*)]
    \item Montrer que $F$ est un sous-espace vectoriel de $\mathbb{R}^3$.
    \item Montrer que $G\subset F$.
    \item Déterminer un vecteur de $F$ n’appartenant pas à $G$.
    \item Décrire géométriquement $F$ et $G$.
\end{enumerate}

\end{document}